\chapter{DETAILED SYSTEM IMPLEMENTATION}
\graphicspath{{Chapter4/}}

In this chapter, we describe the specific algorithmic instantiation of the proposed Asymmetric Multi-Probe LSH SE scheme. Specifically, the scheme is formalized by four main polynomial-time algorithms: \textbf{Setup}, \textbf{BuildIndex}, \textbf{Trapdoor}, and \textbf{Search}.

\section{System Setup (Setup)}
The Setup phase is executed by the Data Owner (DO) to initialize the cryptographic parameters of the Inner Product Predicate Encryption (IPE) framework and the public hashing functions.

\begin{itemize}
    \item \textbf{Input:} A security parameter $\lambda$, and the desired Bloom filter length $m$.
    \item \textbf{Operations:}
    \begin{enumerate}
        \item The DO initializes the bilinear pairing parameters $(\mathbb{G}, \mathbb{G}_T, p, e, g)$ where $p$ is a prime based on $\lambda$ \cite{shen2009predicate}.
        \item The DO generates the Master Secret Key ($MSK$) and the corresponding Public Key ($PK$) required for the IPE scheme over vectors of length $m$. 
        \item The DO selects $l$ independent LSH functions from the $p$-stable LSH family $\mathcal{H} = \{h: \{0,1\}^{26^2} \rightarrow \{0,1\}^m\}$ \cite{datar2004locality}.% [cite: 197]
    \end{enumerate}
    \item \textbf{Output:} The DO keeps $MSK$ completely private. The DO publicly broadcasts the system parameters $PARAMS = \{PK, \mathbb{G}, \mathbb{G}_T, e, \mathcal{H}, m, l\}$.
\end{itemize}

\section{Index Generation (BuildIndex)}
Before outsourcing a document, the Data Owner (or an authorized data provider holding the $PK$) must generate its secure index. 

\begin{itemize}
    \item \textbf{Input:} A document $\mathcal{D}$, the public key $PK$, and the hashing parameters.
    \item \textbf{Operations:}
    \begin{enumerate}
        \item \textbf{Keyword Extraction:} Extract the set of distinct keywords $\mathcal{W_D} = \{w_1, w_2, \dots \}$ from the document $\mathcal{D}$.% [cite: 199]
        \item \textbf{Vectorization:} Convert each keyword $w_i$ into its $26^2$-dimensional bigram vector representation \cite{wang2014privacy}.% [cite: 139, 141]
        \item \textbf{Bloom Filter Insertion:} Initialize an $m$-bit Bloom filter $\mathcal{I_D}$ to all zeros. For each bigram vector, apply the $l$ LSH functions $h_j \in \mathcal{H}$ (where $1 \le j \le l$) and set the corresponding bits in $\mathcal{I_D}$ to $1$.% [cite: 198, 200]
        \item \textbf{Asymmetric Encryption:} Treat the final Bloom filter $\mathcal{I_D}$ as an $m$-dimensional vector. Encrypt this vector using the IPE encryption algorithm with the Public Key ($PK$) to produce the ciphertext index $CT_{\mathcal{I}}$ \cite{shen2009predicate}.
    \end{enumerate}
    \item \textbf{Output:} The secure index $CT_{\mathcal{I}}$, which is outsourced to the Cloud Server alongside the symmetrically encrypted document payload.
\end{itemize}

\section{Search Token Generation (Trapdoor)}
When an authorized user wants to search for documents, they cannot generate the search token themselves because they do not possess the $MSK$. They must interact securely with the DO.

\begin{itemize}
    \item \textbf{Input:} The user's query containing fuzzy keywords $\mathcal{Q}_{words}$, and the DO's $MSK$.
    \item \textbf{Operations:}
    \begin{enumerate}
        \item \textbf{Query Formulation:} The user transforms the keywords in $\mathcal{Q}_{words}$ into bigram vectors and maps them into an $m$-bit Bloom filter query vector $\mathcal{Q}$ using the same $l$ public LSH functions $h_j \in \mathcal{H}$.% [cite: 202]
        \item \textbf{Token Request:} The user securely transmits the plaintext vector $\mathcal{Q}$ (or the raw keywords, depending on the trust model) to the Data Owner.
        \item \textbf{Token Generation:} The DO utilizes the Master Secret Key ($MSK$) and the IPE token generation algorithm to encrypt the query vector $\mathcal{Q}$ into a secure search token $TK_{\mathcal{Q}}$ \cite{shen2009predicate}.
    \end{enumerate}
    \item \textbf{Output:} The authorized trapdoor token $TK_{\mathcal{Q}}$, which the user submits to the Cloud Server.
\end{itemize}

\section{Secure Search and MP-LSH Evaluation (Search)}

Upon receiving the token from the user, the Cloud Server evaluates the match against the stored encrypted indexes without learning the plaintext keywords or index contents.

\begin{itemize}
    \item \textbf{Input:} The user's trapdoor token $TK_{\mathcal{Q}}$, the encrypted document indexes $CT_{\mathcal{I}}$, and the Multi-Probe perturbation parameter $c$.
    \item \textbf{Operations:}
    \begin{enumerate}
        \item \textbf{Primary Evaluation:} For each encrypted index $CT_{\mathcal{I}}$ in the database, the CS executes the IPE public evaluation algorithm $e(CT_{\mathcal{I}}, TK_{\mathcal{Q}})$. Due to the bilinear pairings, this yields the inner product $\langle \mathcal{I_D}, \mathcal{Q} \rangle$, representing the number of matched bits \cite{shen2009predicate}.
        \item \textbf{Multi-Probe Expansion:} To minimize false negatives for heavily misspelled queries, the CS generates systematic perturbation vectors $c = (\delta_1, \dots, \delta_l)$ \cite{lv2007multi}. 
        \item \textbf{Adjacent Bucket Checking:} The CS mathematically shifts the query token's evaluation space by the perturbation vectors $c$ to check adjacent hashing buckets. If an adjacent bucket yields a higher inner product score, the CS dynamically updates the similarity rank.
        \item \textbf{Ranking:} The CS sorts the documents based on the final, probe-adjusted inner product scores. 
    \end{enumerate}
    \item \textbf{Output:} The CS returns the top-$k$ ranked encrypted documents to the querying user.% [cite: 204]
\end{itemize}